\section{MPI scaling}\label{sec:mpi}

\subsection{Set-up}
All runs use Vanda's 72-core nodes (\cref{tab:environment}), one core per rank,
ranks filling the first socket, then the second, then the next node. A job owns its
nodes exclusively. Configurations of up to 72 cores were measured in one-node jobs
and larger ones in a four-node job; the 72-rank configurations were measured in
both and agree within \n{strong.anchor.maxdev.pct}\%. Each configuration warms up
for at least five seconds and is then measured three times; the time of a repetition
is the maximum over ranks of the attention phase (state reset, communication,
computation, final division), and the reported time is the median of the three.
Data generation, window creation and validation are outside the timed phase. The
three repetitions follow each other inside one job, so their spread describes
short-term fluctuation, not the variation between allocations. For the
\n{rep.long.n} timed CPU configurations that last at least 50\,ms, they differ by
\n{rep.long.median.pct}\% in the median and by less than \n{rep.long.p90.pct}\% for
nine in ten; \n{rep.long.above10.pct}\% differ by more than 10\%, all of them in the
four-node job. Configurations shorter than 50\,ms are noisier (median
\n{rep.short.median.pct}\%), which concerns the smallest shards of
\cref{fig:smallc}. Differences of one or two per cent between two configurations
are therefore not resolved by this protocol, and I do not interpret them.
Unless stated otherwise $d=d_v=64$, one head, no mask, and the tile sizes of
\cref{sec:serial}.

\paragraph{The serial baseline.} $T_S$ is the time of the tuned kernel on one core
(S1), measured in the same job and on the same node as the parallel runs it is
compared with. The same loops can also be reached as the single rank of the
no-communication control (M0) or of the ring (M2). Three alternating rounds of the
three give \n{baseline.S1}, \n{baseline.local} and \n{baseline.ring}\,s at
$N=32{,}768$, within \n{baseline.gap.pct}\% of each other: the MPI driver adds
nothing measurable to a single rank, so a speedup below is a statement about
parallel execution and not about two different serial codes.

\subsection{Strong scaling}
\inputifexists{report/generated/table_strong.tex}

\begin{figure}[t]
\centering
\figureifexists{strong_scaling.pdf}
\caption{Strong scaling to 288 cores. (a) Speedup of the overlapped ring for three
problem sizes. (b) Parallel efficiency at the largest size of the five variants in
the strong-scaling grid (M1 and M7 from eight ranks; M1b and M8 appear in
\cref{fig:smallc}). (c) The efficiency of M3 in (b) written as the product of two
factors: the efficiency of the control M0, and the speed of M3 relative to M0.}
\label{fig:strong}
\end{figure}

\Cref{tab:strong,fig:strong} give the main result. For $N=\n{strong.Nbig.h}$ the
overlapped ring reaches a speedup of \n{strong.S.big.72} on one node (72 cores,
efficiency \n{strong.E.big.72}) and \n{strong.S.big.288} on four nodes (288 cores,
efficiency \n{strong.E.big.288}); the serial run of \n{strong.TS.big}\,s becomes
\n{strong.T.big.288}\,s. For the smallest size, $N=\n{strong.Nsmall.h}$, the
efficiency at 288 cores is \n{strong.E.small.288}. The variants agree inside one
node and separate beyond it (\cref{fig:strong}b): at 288 cores the blocking ring
reaches \n{strong.E.ring.big.288}, the all-gather \n{strong.E.allgather.big.288} and
the one-sided pull \n{strong.E.rma_get.big.288}. Up to eight ranks, where the
processor runs at the same clock as in the serial run, every measured variant and
size lies between \n{strong.Elow.min} and \n{strong.Elow.max}: the decomposition
itself costs at most \n{strong.Elow.loss.pct}\%. Two questions follow: why does the
efficiency fall inside a single node, where communication is cheap, and why does it
fall further for small problems and across nodes?

\subsection{A loaded node is slower: the control M0}\label{sec:hardware}
The control M0 answers the first question, because its timed phase contains no MPI
call at all. With one rank a core delivers \n{hw.gflops.1}\,GFLOP/s; with all 72
cores of the node busy each delivers \n{hw.gflops.72}\,GFLOP/s, a factor of
\n{hw.eff.72} (\cref{fig:hardware}a). The processor's clock, sampled during the
same runs, moves the same way: \n{hw.freq.9}\,GHz with up to nine active cores of a
socket, \n{hw.freq.18}\,GHz with eighteen, \n{hw.freq.36}\,GHz with the socket full
and \n{hw.freq.72}\,GHz with both sockets full (\cref{fig:hardware}b), a ratio of
\n{hw.freq.ratio} that accounts for most of the loss. Placing 36 ranks as 18 per
socket instead of 36 on one socket keeps the clock at \n{hw.freq.spread36}\,GHz and
is \n{hw.spread.pct}\% faster, so the limit acts per socket: this is a processor
whose turbo clock depends on the number of active cores. The 256-bit build
declines by a similar factor (\n{hw.w256.eff.72}) from a lower level
(\n{hw.w256.gflops.1} to \n{hw.w256.gflops.72}\,GFLOP/s per core), so the wide
vectors are not what triggers it.

The sampled frequency is an instantaneous reading, so a second job measured the
two things directly. It ran the control with every rank under \code{perf stat}, for
a fixed amount of work that lasts \n{clock.seconds.min} to \n{clock.seconds.max}\,s
per rank count ($N=\n{clock.N.h}$). Cycles divided by CPU time give the clock a
rank actually ran at, averaged over its run; cycles divided by query--key pairs
give the work a cycle does (\cref{fig:hardware}b,c). One rank runs at
\n{clock.ghz.1}\,GHz (the average includes the first seconds at the base clock) and
needs \n{clock.cpp.1} cycles per pair. Up to 18 ranks neither number moves. With one
socket full the clock is \n{clock.ghz.36}\,GHz and with both full
\n{clock.ghz.72}\,GHz, a factor of \n{clock.fratio.72}; the cycles per pair rise to
\n{clock.cpp.36} at 36 ranks and \n{clock.cpp.72} at 72, a factor of
\n{clock.wratio.72} in work per cycle. The efficiency of the control at 72 ranks,
\n{clock.timed.72} in this run, is thus the product of a clock factor of
\n{clock.fratio.72} and a work-per-cycle factor of \n{clock.wratio.72}; for all
eight configurations the product of the two counter-based factors agrees with the
timed efficiency within \n{clock.product.gap.pct}\%. The 36 ranks spread over both
sockets run at \n{clock.ghz.spread36}\,GHz with the same work per cycle as the 36
packed ones (\n{clock.wratio.spread36} against \n{clock.wratio.36}): the
difference that the counters show between the two placements is the clock. Power
and temperature were not measured, so
\emph{why} the processor lowers its clock is not established here; that it does,
and by how much, is.

\begin{figure}[t]
\centering
\figureifexists{hardware.pdf}
\caption{The no-communication control M0 on one node ($N=65{,}536$). (a) Useful
throughput per core against the number of active cores. (b) Clock: frequency of the
busiest cores sampled every 0.25\,s during the measured repetitions (coloured), and
the average over a whole run from hardware counters, cycles per CPU second, in a
separate job on another node (black). (c) Query--key pairs per cycle from the same
counters, relative to one rank.}
\label{fig:hardware}
\end{figure}

I therefore write the efficiency as a product,
\[
 E=\frac{T_S}{P\,T_P}=\underbrace{\frac{T_S}{P\,T_{\text{M0}}}}_{E_0}\times
 \underbrace{\frac{T_{\text{M0}}}{T_P}}_{E_\text{comm}} .
\]
This is an identity, not a measurement of two causes: $E_0$ is the efficiency of
the control and $E_\text{comm}$ is the speed of a variant \emph{relative to that
control}. Reading $E_\text{comm}$ as ``the cost of communication'' needs the control
to differ from the variant in nothing else, which holds only approximately. M0 runs
the same kernel on the same blocks in the same order, but it holds all $N$ rows of
keys and values and copies each block out of that array, where a ring receives the
block into its work buffer; memory footprint, one copy per step and cache contents
differ. One check supports the reading: the time that the ranks spend inside the
kernel is the same in both, within \n{strong.kernel.ratio.min} to
\n{strong.kernel.ratio.max} over the \n{strong.kernel.ratio.n} strong-scaling
configurations of the two rings (\n{strong.kernel.ratio.big.288} at 288 ranks for
the largest problem), so what $E_\text{comm}$ loses is spent outside the kernel.
\Cref{fig:strong}c shows both factors for the largest problem: at 72 cores
$E_0=\n{strong.Ehw.big.72}$ and $E_\text{comm}=\n{strong.Ecomm.big.72}$; at 288 cores
$E_0=\n{strong.Ehw.big.288}$ and $E_\text{comm}=\n{strong.Ecomm.big.288}$. Most of the
loss is already present in the control, on one node and on four.

$E_0$ is therefore a product of two measured factors. The first is the clock. The
second, the work per cycle, behaves like a cost of the granularity of the
decomposition: every rank transposes all $N$ keys into tiles for its own $c$
queries, and the last query tile of a rank and the last key tile of every shard
are partial, so the same pairs cost a little more when they are cut into more and
smaller pieces. The counters are consistent with this reading, because the factor
follows the rows per rank more than the load. At 1{,}820 rows per rank it is \n{clock.wratio.36} with 36 packed ranks and
\n{clock.wratio.spread36} with 36 spread ranks, and the strong-scaling run with 72
ranks at $N=131{,}072$ implies \n{clock.implied.w.big} ($E_0$ divided by the clock
factor); at 910 rows it is \n{clock.wratio.72} from the counters and
\n{clock.implied.w.mid} implied by the strong-scaling run; at 455 rows the
strong-scaling run implies \n{clock.implied.w.small}. That is why $E_0$ at 72 ranks
falls from \n{strong.E0.big.72} for $N=131{,}072$ to \n{strong.E0.small.72} for
$N=32{,}768$, and why it is lower at 288 ranks than at 72 for the same $N$ although
all four nodes run at the same clock. Granularity is not the only possible
cause, though. Contention for caches and memory when both sockets are full would
also lower the work per cycle, and the 72-rank points cannot tell the two apart;
the agreement at equal rows per rank under different loads speaks for
granularity as the main part, without excluding a contribution of contention. The nodes themselves are equally fast: with all 288 cores busy the
median compute time of a rank differs by \n{nodes.spread.pct}\% between the four
nodes, and the fastest and the slowest of the 288 ranks differ by
\n{nodes.rank.spread.pct}\%.

\subsection{Speed relative to the control against rows per rank}\label{sec:collapse}
\begin{figure}[t]
\centering
\figureifexists{collapse.pdf}
\caption{$E_\text{comm}=T_\text{M0}/T$ of the two-sided variants against rows per
rank, for the strong-scaling, weak-scaling and small-shard runs (one thread per
rank). Open markers: all ranks on one node; filled: several nodes. Dashed:
\cref{eq:model} with parameters from the two-rank microbenchmark. Solid: the same
form with one constant fitted to the rings with \n{model.thin.ranks} single-thread
ranks on one node and $c\ge\n{model.cmin.fit}$ (\n{model.fit.n} points); all other
points took no part in the fit. Dotted: no fitted constant, the compute time of a
block divided by itself plus the exchange time measured with 72 ranks and no
computation.}
\label{fig:collapse}
\end{figure}

The second question is approached with \cref{eq:model}: for a given configuration
communication should matter through $c=N/P$. \Cref{fig:collapse} plots
$E_\text{comm}$ of every two-sided run with one thread per rank against $c$. Inside
one node, runs with different $N$ and $P$ fall close to one curve. With thin ranks
the ring keeps \n{collapse.E.c512.ring} of the control's speed at 512 rows per
rank, \n{collapse.E.c128.ring} at 128 and \n{collapse.E.c16.ring} at 16. (At $c=8$ a
repetition lasts a third of a millisecond; those points scatter and I draw no
conclusion from them.)

\paragraph{How well does the model predict the curve?} The dashed line uses
parameters measured independently of the scaling runs: a per-exchange cost
$\alpha=\n{model.alpha.us}\,\mu$s and a bandwidth $B=\n{model.B.gbs}$\,GB/s for
megabyte messages, both from a ping-pong between two ranks of an otherwise idle
node, and $F=\n{model.F.gflops}$\,GFLOP/s per core under full load. They give
$c^*=\n{model.cstar}$ and $c_\alpha=\n{model.calpha}$ rows. The prediction has the
right shape but is clearly too optimistic: at $c=128$ it gives \n{model.pred.c128}
where the rings measure \n{collapse.E.c128.ring} and \n{collapse.E.c128.overlap}.
The measured losses nevertheless have the form of the model's bandwidth term. The
\n{model.fit.n} points used for the fit ($\n{model.cmin.fit}\le c\le\n{model.cmax.fit}$,
\n{model.thin.ranks} ranks on one node) follow $E_\text{comm}=1/(1+c^*/c)$ with a
single constant, $c^*=\n{model.cstar.fit}$ rows, to within \n{model.fit.maxdev}
(solid line); the latency term is not needed. The fitted constant is
\n{model.cstar.ratio} times the predicted one. That is a statement about this
constant, not a uniform factor on every communication cost.

\paragraph{Why the microbenchmark was too optimistic.} The ping-pong moves one
message one way between two ranks of an idle node. A ring step is different in two
respects: every rank sends \emph{and} receives, and all ranks do so at once. A
second microbenchmark measures exactly that: the two \code{Sendrecv} calls of a
ring step, repeated with nothing computed in between, for \n{ringx.ranks.list}
ranks of one node. For 128-row shards one step takes \n{ringx.us.P2.c128}\,$\mu$s
between two ranks, \n{ringx.pair.over.model.c128} times what the ping-pong
constants give, and \n{ringx.us.P72.c128}\,$\mu$s when 72 ranks exchange at the same
time, \n{ringx.crowd.c128} times more again. A straight line through the 72-rank
times for $16\le c\le512$ gives a bandwidth of \n{ringx.B.gbs}\,GB/s where the
ping-pong gave \n{model.B.gbs}, and a fixed cost of \n{ringx.alpha.us}\,$\mu$s per
step. With that bandwidth \cref{eq:model} gives $c^*=\n{ringx.cstar}$ rows, against
\n{model.cstar.fit} fitted to the scaling runs and \n{model.cstar} from the
ping-pong. Calibrating the constant in the situation in which the ring actually
runs therefore explains most of the difference: the larger part of the factor of
\n{model.cstar.ratio} is the crowd, 72 ranks copying through the same memory system
at once, and the rest is the two-way exchange. How far the model's form itself
holds can be checked by using the measured exchange time directly, without any
constant: the time of one block of the control divided by that time plus the
exchange time (dotted in \cref{fig:collapse}). It reproduces the measured ring at
256 and 512 rows per rank (\n{ringx.pred.E.c256} and \n{ringx.pred.E.c512} against
\n{ringx.meas.E.c256} and \n{ringx.meas.E.c512}) and is too pessimistic for smaller
shards (\n{ringx.pred.E.c128} against \n{ringx.meas.E.c128} at 128 rows,
\n{ringx.pred.E.c64} against \n{ringx.meas.E.c64} at 64). So the simple form,
compute time plus an exchange time that does not depend on the computation, is
confirmed for shards of 256 rows and more and not below. A plausible reason for
the deviation is that in a real run the ranks do not all enter the exchange at the
same instant, as they do in the tight loop of the microbenchmark; I did not test
this.

\paragraph{How far does the fitted curve carry?} The runs that were not used for
the fit test it. The \n{model.heldout.one.n} one-node configurations of the two
rings in the strong- and weak-scaling runs, with 2 to 72 ranks, lie within
\n{model.heldout.one.maxdev} of the curve. Across nodes the overlapped ring stays
within \n{model.heldout.cross.overlap.maxdev} (\n{model.heldout.cross.overlap.n}
configurations), while the blocking ring falls up to
\n{model.heldout.cross.ring.below} and the all-gather up to
\n{model.heldout.cross.allgather.below} below it; at 288 ranks and $N=131{,}072$,
for example, $E_\text{comm}$ is \n{strong.Ecomm.ring.big.288} and
\n{strong.Ecomm.allgather.big.288} against \n{strong.Ecomm.big.288}. Rows per rank
therefore organise the two-sided results well, but they are not the only variable.
Threads per rank (below), the network between the ranks, the blocking or
overlapped schedule, and for the one-sided pull even the order of the fetches
(\cref{sec:progress}) change the cost at equal $c$. The fit also yields a rule of
thumb for thin ranks inside a node, half the speed at $c=c^*\approx\n{model.cstar.fit}$
rows and a 5\% loss at $19\,c^*\approx\n{model.c95.fit}$ rows; the second number is an
extrapolation of the fitted form, not a measured threshold.

\paragraph{One-sided inside a node.} The one-sided pull does not follow this curve.
With the same thin ranks it stays between \n{collapse.get.min} and
\n{collapse.get.max} of the control's speed down to 16 rows per rank. A plausible
reason is that a Get from a window in shared memory can be served by the reader
alone with a single copy, whereas a matched send and receive needs both ranks;
how many copies each path makes depends on the library's shared-memory transport
and the message size, which I did not inspect. \Cref{sec:progress} shows that this
advantage turns into its opposite between nodes.

\paragraph{Faster ranks are harder to feed.} The second prediction of
\cref{sec:model} can be tested inside the node as well. With \n{fat.ranks} ranks of
\n{fat.threads} threads the kernel of a rank is many times faster, while the
exchange of a shard of given size does not speed up in the same proportion. At 2{,}304 rows per rank the ring
then keeps \n{fat.E.ring.T18.c2304} and the overlapped ring
\n{fat.E.overlap.T18.c2304}, where the thin-rank curve gives \n{fat.thin.c2304}; the
loss falls below a few per cent only near 18{,}000 rows
(\n{fat.E.ring.T18.c18432}).

\begin{figure}[t]
\centering
\figureifexists{small_c.pdf}
\caption{Run time relative to the control M0 when ranks own few rows. (a) 288
single-thread ranks on four nodes. (b) 16 ranks of 18 threads on the same four
nodes; open markers: the repeat of three sizes in a later job, two rounds.
Logarithmic vertical axes.}
\label{fig:smallc}
\end{figure}

\Cref{fig:smallc} shows the regime of few rows per rank across four nodes for every
variant. With 288 thin ranks and 128 rows each, the overlapped ring takes
\n{smallc.thin.overlap.c128} times as long as the control, the blocking ring
\n{smallc.thin.ring.c128} and the all-gather \n{smallc.thin.allgather.c128}; at 8
rows the factors are \n{smallc.thin.overlap.c8}, \n{smallc.thin.ring.c8} and
\n{smallc.thin.allgather.c8}. The nonblocking all-gather is much worse
(\n{smallc.thin.iallgather.c128} and \n{smallc.thin.iallgather.c8}), and the
one-sided variants are in a different class: the push with a fence per step takes
\n{smallc.thin.rma_put.c128} and \n{smallc.thin.rma_put.c8} times the control, the
pull \n{smallc.thin.rma_get.c128} and \n{smallc.thin.rma_get.c8} times. With 16 ranks
of 18 threads (panel b) a rank computes much faster, and the same ordering appears
at roughly ten times more rows per rank: at 1{,}152 rows the overlapped ring needs
\n{smallc.fat.overlap.c1152}, the blocking ring \n{smallc.fat.ring.c1152} and the
pull \n{smallc.fat.rma_get.c1152} times the control. One point of panel (b) is out
of line: the overlapped ring at 4{,}608 rows took \n{smallc.fat.overlap.c4608} times
the control in all three repetitions of that run, with the compute time itself
inflated, while the same configuration with polling, measured minutes later, took
\n{smallc.fat.overlappoll.c4608} times. That could be a disturbance of the run or a
property of the nonblocking exchange at this size that polling cures, and three
consecutive repetitions cannot tell the two apart. I therefore measured the
configuration again in a separate job on the same four nodes, together with the
neighbouring sizes, in two rounds (open markers). The overlapped ring took
\n{fatcheck.overlap.c4608.r1} and \n{fatcheck.overlap.c4608.r2} times the control at
4{,}608 rows (blocking ring: \n{fatcheck.ring.c4608.r1} and
\n{fatcheck.ring.c4608.r2}; with polling: \n{fatcheck.overlappoll.c4608.r1} and
\n{fatcheck.overlappoll.c4608.r2}), in line with its neighbours. The first
measurement was a disturbed run, and polling did not help. The disturbed point
stays in the figure, because it is what that job measured.

Two further settings of the one-sided pull were tested at 288 ranks. If every rank
fetches the shards in the same ascending order, so that all ranks ask the same
owner at the same time, the run takes \n{order.ratio} times as long as with the
staggered order used everywhere else (\n{order.T.ascending}\,s against
\n{order.T.staggered}\,s). Equal rows per rank, different order, a factor of \n{order.ratio}: for the pull, $c$
is clearly not the only variable. Calling \code{MPI\_Testall} between batches of
query tiles to give the library a chance to progress does not help with thin
ranks: the speed changes by a factor between \n{poll.worst.overlap} and
\n{poll.best.overlap} for the overlapped ring and between \n{poll.worst.rma_get} and
\n{poll.best.rma_get} for the pull.

\subsection{Progress: when does a transfer complete while the peer computes?}\label{sec:progress}
\begin{figure}[t]
\centering
\figureifexists{netbench.pdf}
\caption{MPI microbenchmark with two ranks. (a) Time per message against size for
two-sided ping-pong and one-sided Get, inside a node and between nodes. (b) Time
until a transfer completes while the \emph{other} rank computes for 0.2\,s without
calling MPI: a receive from a busy sender, and a Get from a busy target. The rank
that waits is inside MPI the whole time.}
\label{fig:netbench}
\end{figure}

M3 and M7 post their transfers before computing, hoping that the bytes move in the
meantime. Whether they do is a property of the MPI library and the network, which
\cref{fig:netbench} probes for one situation: one rank waits inside MPI for a
transfer while its peer computes and does not call MPI. Inside a node (shared
memory) a small message costs \n{net.lat.intra.us}\,$\mu$s, a megabyte moves at
\n{net.bw.intra.gbs}\,GB/s, and for the sizes tested (64\,KiB to 16\,MiB) the
transfer completes within its normal time although the peer is busy. Between nodes
the picture is different. The Open MPI installation selects TCP over Ethernet: a
small message costs \n{net.lat.inter.us}\,$\mu$s, a megabyte moves at
\n{net.bw.inter.gbs}\,GB/s, and the transfer completes only after the busy rank
returns to MPI, \n{net.busy.ms}\,ms later, exactly the time it computed. The same
holds for the one-sided Get, which this installation carries out in software on
top of TCP and which needs the target to call MPI. In this test, then, a
``nonblocking'' or ``one-sided'' transfer across nodes does not advance while the
peer computes. The test does not cover every case of the application: there both
ranks compute at once, and a transfer can still advance whenever both happen to
be inside MPI. It does say that posting a transfer early is not enough on this
system, and the following measurements of the application agree with it.

\begin{figure}[t]
\centering
\figureifexists{traces.pdf}
\caption{One rank of eight (two per node on four nodes, nine threads each,
$N=32{,}768$): time per ring step, median of the three repetitions, on logarithmic
axes. (a) Time blocked in an MPI wait. (b) Time in the compute kernel. These are
program events (when a call started and returned), not a recording of the bytes
on the wire. The last step has no exchange.}
\label{fig:traces}
\end{figure}

\Cref{fig:traces} shows the consequence step by step for one rank of eight spread
over four nodes. The blocking ring computes for \n{trace.compute.ring}\,s and is
blocked in MPI for \n{trace.wait.ms.ring}\,ms (\n{trace.wait.pct.ring}\% of the
repetition, about \n{trace.wait.step.ms.ring}\,ms per step); the overlapped ring is
blocked for \n{trace.wait.ms.overlap}\,ms (\n{trace.wait.pct.overlap}\%). The
difference is small in absolute terms: the two runs take \n{trace.total.ring} and
\n{trace.total.overlap}\,s, \n{trace.total.gain.pct}\% apart. Given the
microbenchmark, it is more plausibly a difference of schedule than of overlap: M3
posts its four requests together, one step ahead, and waits for all of them at
once, whereas M2 performs two blocking exchanges one after the other when it has
finished computing. The one-sided pull is blocked for \n{trace.wait.ms.rma_get}\,ms
(\n{trace.wait.pct.rma_get}\%), unevenly from step to step, and the run takes
\n{trace.total.rma_get}\,s instead of \n{trace.total.overlap}\,s. This is what one
expects if each Get waits until its target, which is busy with its own block, next
enters MPI, and it is consistent with the collapse of M7 in \cref{fig:strong}b,
where it is \n{strong.slow.rma_get.big.288} times slower than the control at 288
ranks for $N=131{,}072$ and \n{strong.slow.rma_get.small.288} times for $N=32{,}768$.
The traces show when calls return; they do not by themselves separate waiting for
the target from congestion or from the cost of the software protocol.

\paragraph{Placement can outweigh the network.} The same 72 ranks ($N=65{,}536$)
were run packed on one node and spread over four nodes, 18 per node. The spread
run is \n{placement.gain.pct.overlap}\% faster (\n{placement.spread.overlap}\,s
against \n{placement.packed.overlap}\,s for the overlapped ring, and the same
\n{placement.gain.pct.local}\% for the control), although every exchange now crosses
the Ethernet: with 18 active cores per node the processors run at
\n{hw.freq.18}\,GHz instead of \n{hw.freq.72}\,GHz, and at 910 rows per rank the
ring's communication costs less than the clock gains.

\subsection{Weak scaling}
\begin{figure}[t]
\centering
\figureifexists{weak_scaling.pdf}
\caption{Weak scaling, logarithmic axes. (a) Rows per rank fixed ($N=1024P$, up to
288 ranks): the work per rank grows with $P$. (b) Work per rank fixed
($N=16{,}384\sqrt P$, up to 256 ranks).}
\label{fig:weak}
\end{figure}

For a quadratic algorithm ``weak scaling'' has two meanings and I measured both
(\cref{fig:weak}). If the \emph{rows per rank} are kept constant ($N\propto P$,
constant memory per rank), every rank must still meet all $N$ keys, so its work
grows in proportion to $P$ and the run time must too; the model predicts constant
\emph{efficiency}, not constant time. The measured time follows the $\propto P$
line from 1 to 288 ranks. In terms of efficiency, $P\,T_1/T_P$ is
\n{weak.rows1024.E.overlap.72} at 72 ranks and \n{weak.rows1024.E.overlap.288} at
288 for the overlapped ring, and \n{weak.rows1024.E.local.72} and
\n{weak.rows1024.E.local.288} for the control: constant, as predicted, and again
already present in the control. The corresponding scaled speedup (the time one
core would need for the $P$ times longer sequence, $P^2T_1$, over the measured
time) is \n{weak.rows1024.S.overlap.72} at 72 ranks and
\n{weak.rows1024.S.overlap.288} at 288. The blocking ring reaches
\n{weak.rows1024.E.ring.288} at 288 ranks and the one-sided pull
\n{weak.rows1024.E.rma_get.288}. If instead the \emph{work per rank} is
kept constant ($N\propto\sqrt P$), the ideal time is flat. The measured time is
\n{weak.work16384.T1}\,s on one core and \n{weak.work16384.T.overlap.16}\,s on 16,
rises to \n{weak.work16384.T.overlap.64}\,s on 64 cores, where the node is nearly
full and the clock has dropped, and stays at \n{weak.work16384.T.overlap.256}\,s on
256 cores and four nodes (efficiency $T_1/T_P=\n{weak.work16384.E.overlap.256}$, scaled
speedup $P\,T_1/T_P=\n{weak.work16384.S.overlap.256}$; control
\n{weak.work16384.E.local.256}). Going from one node to four costs nothing further
at this size, because a rank still owns about 1{,}000 rows.

\subsection{Causal masking: balance against traffic}\label{sec:causal}
\begin{figure}[t]
\centering
\figureifexists{concept_causal.pdf}
\caption{Why a causal mask unbalances a contiguous partition. Allowed pairs
(coloured by the rank that owns the query) for $N=32$, $P=4$. (a) Contiguous: the
last rank owns the long rows. (b) Zigzag: every rank owns one early and one late
chunk. (c) Resulting share of the work.}
\label{fig:causalconcept}
\end{figure}

With a causal mask query $i$ uses $i+1$ keys, so the work is a triangle
(\cref{fig:causalconcept}). A contiguous partition gives the last rank
$(2P-1)/P^2$ of all pairs, almost twice the average, and the first rank almost
nothing. The run then takes as long as the unmasked one although half the work is
gone. The zigzag partition pairs chunk $r$ with chunk $2P-1-r$, whose row lengths
add up to the same total for every $r$, so each rank gets exactly $1/P$ of the
pairs.

\begin{figure}[t]
\centering
\figureifexists{causal.pdf}
\caption{Causal mask, $N=131{,}072$. (a) Compute time of each of eight ranks
(logarithmic axis). (b) Run time relative to the unmasked run of the overlapped
ring (8 and 72 ranks on one node, 288 on four; logarithmic axis). (c) Bytes
received by the one-sided pull relative to the ring, summed over ranks.}
\label{fig:causal}
\end{figure}

\Cref{fig:causal} confirms the picture. With eight ranks and a contiguous partition
the compute time per rank runs from \n{causal.rank.min.contiguous}\,s to
\n{causal.rank.max.contiguous}\,s, a factor of \n{causal.rank.spread.contiguous};
with the zigzag partition every rank needs between \n{causal.rank.min.zigzag} and
\n{causal.rank.max.zigzag}\,s. The run time follows the slowest rank: relative to the
unmasked run the contiguous causal run takes \n{causal.ratio.overlap.contiguous.8}
with 8 ranks and \n{causal.ratio.overlap.contiguous.72} with 72, so that halving
the work saves almost nothing, while the zigzag run takes
\n{causal.ratio.overlap.zigzag.8} and \n{causal.ratio.overlap.zigzag.72}.
On four nodes with 288 ranks the result is the same, \n{causal.ratio.overlap.contiguous.288}
for the contiguous and \n{causal.ratio.overlap.zigzag.288} for the zigzag partition.

The two partitions also differ in traffic. A ring must pass every shard to every
rank whatever the mask. The one-sided pull fetches only shards that contain at
least one allowed key: with the contiguous partition rank $r$ needs the shards of
ranks $0,\dots,r$ only, and the pull moves \n{causal.bytes.contiguous.8} of the
ring's bytes at 8 ranks and \n{causal.bytes.contiguous.288} at 288; with the zigzag
partition every rank owns late rows and needs everything
(\n{causal.bytes.zigzag.288}). Balance and traffic therefore pull in opposite
directions, and which one wins depends on what a fetch costs. Inside a node, where
it is cheap, balance wins for the pull as well: zigzag is
\n{causal.zigzag.gain.rma_get.72} times faster than contiguous with 72 ranks. On
four nodes, where every fetch waits for its target (\cref{sec:progress}), the
order reverses: the pull takes \n{causal.T.rma_get.contiguous.288}\,s with the
contiguous partition and \n{causal.T.rma_get.zigzag.288}\,s with zigzag,
\n{causal.zigzag.cost.pct.rma_get.288}\% more, because halving the traffic now saves
more than balancing the work. Neither is competitive there: the ring with the
zigzag partition needs \n{causal.T.overlap.zigzag.288}\,s. The bytes confirm the
traffic side of this account; the times alone do not separate waiting, congestion
and protocol cost. The choice of partition is thus conditional: zigzag for a ring,
and for a pull only while fetches are cheap.

\subsection{Memory}
\begin{figure}[!htb]
\centering
\figureifexists{memory.pdf}
\caption{Peak resident memory of the largest rank, 72 ranks on one node, measured
in separate runs without validation. Dotted: the one-sided variant with its shared
window counted once instead of once per rank.}
\label{fig:memory}
\end{figure}

\Cref{fig:memory} shows the price of the all-gather. At $N=\n{mem.N.h}$ on
\n{mem.ranks} ranks, a rank holds \n{mem.allgather.mb}\,MiB with the all-gather (the
\n{mem.kv.mb}\,MiB of the complete $K$ and $V$ plus its own shard and state) and
\n{mem.ring.mb}\,MiB with either ring, a factor of \n{mem.ratio}; for the node this
is \n{mem.node.allgather.gb}\,GiB against \n{mem.node.ring.gb}\,GiB. The gap grows
with $P$, because the all-gather stores $N$ rows per rank and the ring $2N/P$. The
one-sided variant also shows \n{mem.rma_get.mb}\,MiB resident, but that is an
accounting effect: its shards live in a shared-memory window that every rank of the
node maps, and the operating system's resident count attributes the whole window to
each of them. Counting every shared page once (the proportional set size) gives
\n{mem.rma_get.pss.mb}\,MiB per rank, close to the ring.
